% !TEX program = pdflatex
\documentclass[11pt]{article}
\usepackage[a4paper,margin=1in]{geometry}
\usepackage{amsmath,amsthm,amssymb,mathtools}
\usepackage{bm,bbm}
\usepackage{microtype}
\usepackage{enumitem}
\usepackage{hyperref}
\usepackage[nameinlink,capitalize]{cleveref}
\hypersetup{colorlinks=true,linkcolor=blue,citecolor=magenta,urlcolor=blue}

% Theorem environments
\theoremstyle{plain}
\newtheorem{theorem}{Theorem}[section]
\newtheorem{lemma}[theorem]{Lemma}
\newtheorem{proposition}[theorem]{Proposition}
\newtheorem{corollary}[theorem]{Corollary}
\theoremstyle{definition}
\newtheorem{definition}[theorem]{Definition}
\theoremstyle{remark}
\newtheorem{remark}[theorem]{Remark}

% Math macros
\newcommand{\R}{\mathbb{R}}
\newcommand{\N}{\mathbb{N}}
\newcommand{\Hcal}{\mathcal{H}}
\newcommand{\Dcal}{\mathcal{D}}
\newcommand{\Qcal}{\mathcal{Q}}
\DeclareMathOperator{\spec}{spec}

\title{A Prime-Indexed Quadratic Model: Operator-Theoretic Foundations, Uniform Excitation Gap, and Reflection Positivity}
\author{Ryan O. Van Gelder, Luis Morat´o de Dalmases, \& Tyler Van Osdol}
\date{\today}

\begin{document}
\maketitle

\begin{abstract}
We formalize a prime-indexed quadratic field model on the direct-sum Hilbert space $\Hcal=\bigoplus_{n\in\N} L^2(\R)$, with instantaneous one-mode Hamiltonians
$$
H_n(t)=\tfrac12\Pi_n^2+\tfrac12\Omega_n(t)^2\,\Phi_n^2,\qquad \Omega_n(t)^2=\Lambda_m^2\omega_n^2+T_{nn}(t),\quad \omega_n=\tfrac{2\pi}{\tau}\ln p_n.
$$
Under a minimal uniform lower bound $T_{nn}(t)\ge-\varepsilon_n$ (with $\varepsilon_n<\Lambda_m^2\omega_n^2$), we prove: (i) mode-wise essential self-adjointness and lower spectral bounds; (ii) a time-uniform, strictly positive single-quantum (excitation) gap $\Delta(t)\ge \sqrt{\delta_*}$ for the normal-ordered many-body Hamiltonian, where $\delta_*:=\inf_n(\Lambda_m^2\omega_n^2-\varepsilon_n)$; (iii) Osterwalder--Schrader reflection positivity of the Euclidean Gaussian measure for even-in-time mass profiles $\Omega_n(t)$. We also state a safe generalization with off-diagonal mode couplings via positive-type kernels. The results are self-contained and intended as drop-in, referee-ready components for a broader Yang--Mills program.
\end{abstract}

\tableofcontents

\section{Introduction}
This article packages a set of operator-theoretic results for a prime-indexed quadratic model that emerged from a larger program aimed at understanding spectral gaps in gauge theories. Our goal here is modest and precise: isolate and prove the functional-analytic statements that are already sufficient to guarantee self-adjointness, uniform-in-time lower bounds, an excitation gap in the many-body setting, and Osterwalder--Schrader (OS) reflection positivity under simple symmetry assumptions.

\paragraph{Model summary.}
Fix $\tau,\Lambda_m>0$ and define prime-indexed frequencies $\omega_n=(2\pi/\tau)\ln p_n$. For each mode $n\in\N$ and time $t\in\R$, define
\begin{equation}\label{eq:Omega}
\Omega_n(t)^2=\Lambda_m^2\omega_n^2+T_{nn}(t),
\end{equation}
with measurable, locally bounded $t\mapsto T_{nn}(t)$. The one-mode Hamiltonians are
\begin{equation}\label{eq:Hn}
H_n(t)=\tfrac12\Pi_n^2+\tfrac12\Omega_n(t)^2\,\Phi_n^2,
\end{equation}
acting on $L^2(\R)$, where $\Phi_n$ is multiplication by $x$ and $\Pi_n=-i\partial_x$.

\paragraph{Main contributions.}
Let $\varepsilon_n\in[0,\Lambda_m^2\omega_n^2)$ and set $\delta_n:=\Lambda_m^2\omega_n^2-\varepsilon_n>0$ and $\delta_*:=\inf_n\delta_n$. Under the minimal hypothesis
\begin{equation}\label{eq:A3}
T_{nn}(t)\ge-\varepsilon_n\quad (\text{all }t),
\end{equation}
we prove:
\begin{itemize}[leftmargin=*]
  \item essential self-adjointness and lower spectral bounds for $H_n(t)$, uniform in $t$;
  \item a time-uniform single-quantum gap $\Delta(t)=\inf_n\Omega_n(t)\ge\sqrt{\delta_*}>0$ for the normal-ordered many-body Hamiltonian on the symmetric Fock space $\Gamma_s(\Hcal)$;
  \item OS reflection positivity when the mass profiles are even in time, $\Omega_n(-t)=\Omega_n(t)$; and
  \item a safe extension to off-diagonal couplings via kernels that are even and of positive type (Bochner).
\end{itemize}

\section{Hilbert space, domains, and assumptions}\label{sec:setup}
\paragraph{Hilbert space and core.} Work on $\Hcal=\bigoplus_{n\in\N}L^2(\R_{x_n})$ with algebraic core
\[
\Dcal_0:=\bigoplus_{n}^{\mathrm{alg}} \mathcal S(\R_{x_n})\qquad(\text{finite support in $n$, Schwartz in each component}).
\]
All operators are first defined on $\Dcal_0$.

\paragraph{Canonical fields and commutators.} On each mode space, $\Phi_n$ is multiplication by $x_n$ and $\Pi_n=-i\,\partial_{x_n}$, with $[\Phi_n,\Pi_m]=i\,\delta_{nm}\,\mathbbm 1$ on $\Dcal_0$.

\paragraph{Assumptions.} For each $n$, $t\mapsto T_{nn}(t)$ is measurable and locally bounded, and
\begin{equation}\label{eq:uniformLower}
T_{nn}(t)\ge-\varepsilon_n,\qquad 0\le \varepsilon_n<\Lambda_m^2\omega_n^2.
\end{equation}
Define $\delta_n=\Lambda_m^2\omega_n^2-\varepsilon_n>0$ and $\delta_*:=\inf_n\delta_n$.

\section{Time-dependent dynamics}\label{sec:dynamics}
\begin{proposition}[Kato form theory for $H_n(t)$]\label{prop:Kato}
Under \eqref{eq:uniformLower}, for each $n$ there exists a unique unitary propagator $U_n(t,s)$ on $L^2(\R)$ solving $i\partial_tU_n(t,s)=H_n(t)U_n(t,s)$ with a common, time-independent form domain. For every compact $I\subset\R$ there is $C_I\ge1$ such that
\[
\|H_n(t)^{1/2}U_n(t,s)\psi\|\ \le\ C_I\,\|H_n(s)^{1/2}\psi\|,\qquad t,s\in I.
\]
\end{proposition}
\begin{proof}[Sketch]
The quadratic forms $\mathfrak h_{n,t}[\psi]=\tfrac12\|\partial_x\psi\|^2+\tfrac12\Omega_n(t)^2\|x\psi\|^2$ are uniformly lower-bounded by \eqref{eq:uniformLower} and locally bounded in $t$, hence Kato stability yields the propagator.
\end{proof}
\noindent\emph{References:} T. Kato, \emph{Perturbation Theory for Linear Operators} (Springer); A. Pazy, \emph{Semigroups of Linear Operators and Applications to Partial Differential Equations} (Springer, Chs.~4--5); K.-J. Engel \& R. Nagel, \emph{One-Parameter Semigroups for Linear Evolution Equations} (Springer).

\section{Essential self-adjointness and lower bounds}\label{sec:essSA}
\begin{lemma}[Essential self-adjointness]\label{lem:essSA}
For fixed $n$ and $t$, the operator $H_n(t)=\tfrac12\Pi_n^2+\tfrac12\Omega_n(t)^2\Phi_n^2$ is essentially self-adjoint on $C_c^\infty(\R)$ (hence on $\mathcal S$). Its closure is self-adjoint, bounded below, with compact resolvent.
\end{lemma}
\begin{proof}
Write $H_n(t)= -\tfrac12\partial_x^2 + V_{n,t}(x)$ with $V_{n,t}(x)=\tfrac12\Omega_n(t)^2x^2$. By \eqref{eq:uniformLower}, $V_{n,t}(x)\ge \tfrac12\delta_n x^2\ge c(1+x^2)$ for some $c>0$. Standard criteria (Faris--Lavine/Kato--Rellich for confining quadratic potentials) imply essential self-adjointness and compact resolvent.
\end{proof}

\begin{lemma}[Form inequality and spectral floor]\label{lem:formBound}
As closed forms on $H^1(\R)\cap\{x\psi\in L^2\}$,
\[
\langle\psi,H_n(t)\psi\rangle=\tfrac12\|\partial_x\psi\|^2+\tfrac12\Omega_n(t)^2\|x\psi\|^2\ \ge\ \tfrac12\|\partial_x\psi\|^2+\tfrac12\delta_n\|x\psi\|^2\ =:\ \langle\psi,h(\delta_n)\psi\rangle.
\]
Therefore $H_n(t)\ge h(\delta_n)$ (KLMN), and
\[
\inf\spec H_n(t)\ \ge\ \tfrac12\sqrt{\delta_n},\qquad \Delta_n(t):=\text{(level spacing)}=\Omega_n(t)\ \ge\ \sqrt{\delta_n}.
\]
\end{lemma}
\begin{proof}
Immediate from \eqref{eq:uniformLower} and the min--max principle applied to the harmonic oscillator with frequency $\sqrt{\delta_n}$.
\end{proof}

\begin{corollary}[Uniform excitation gap; direct sum]\label{cor:gap}
Let $\delta_*:=\inf_n\delta_n$. If $\delta_*>0$, then for the decoupled system $H(t)=\bigoplus_n H_n(t)$ the many-body \emph{excitation} gap (after normal ordering) on the symmetric Fock space $\Gamma_s(\Hcal)$ satisfies
\[
\Delta(t):=\inf_{n}\Omega_n(t)\ \ge\ \sqrt{\delta_*}\ >\ 0\qquad(\text{all }t).
\]
\end{corollary}
\begin{remark}[Fock space and normal ordering]
Let $\Gamma_s(\Hcal)=\bigoplus_{N=0}^{\infty}\Hcal^{\otimes_s N}$. Writing $H_{1\mathrm p}(t):=\bigoplus_n H_n(t)$ and $\mathbf H(t):=\mathrm d\Gamma(H_{1\mathrm p}(t))$, define
\[
:\!\mathbf H(t)\!:\ :=\ \mathbf H(t)-\Big(\sum_n\tfrac12\Omega_n(t)\Big)\mathbbm 1.
\]
On the $N$-particle sector, $:\!\mathbf H(t)\!:\\upharpoonright\ \Hcal^{\otimes_s N}=\sum_{j=1}^N H_{1\mathrm p}(t)^{(j)}$. Hence the single-quantum (excitation) gap equals $\inf_n\Omega_n(t)$ and is independent of the normal-ordering shift.
\end{remark}

\section{Reflection positivity (OS) and minimal fixes}\label{sec:OS}
Consider the Euclidean quadratic action
\[
S_E[\Phi]=\tfrac12\sum_{n}\int_{\R}\big((\partial_t\Phi_n)^2+\Omega_n(t)^2\,\Phi_n^2\big)\,dt,\qquad \Omega_n(t)^2=\Lambda_m^2\omega_n^2+T_{nn}(t).
\]
Let $(\theta f)(t):=f(-t)$. The Gaussian measure with covariance $C_n$ solving $(-\partial_t^2+\Omega_n(t)^2)C_n(\cdot,s)=\delta(\cdot-s)$ is \emph{reflection positive} if, for all test functions $f_n$ supported in $t\ge0$,
\[
\sum_n\iint_{t,s\ge0}\overline{f_n(t)}\,C_n(t,-s)\,f_n(s)\,dt\,ds\ \ge\ 0.
\]

\begin{theorem}[Sufficient condition]\label{thm:OS}
Suppose for each $n$: (i) $\Omega_n\in L^\infty_{\mathrm{loc}}(\R)$, (ii) $\Omega_n(t)\ge \sqrt{\delta_n}>0$, and (iii) \textbf{evenness}: $\Omega_n(-t)=\Omega_n(t)$. Then reflection positivity holds.
\end{theorem}
\begin{proof}[Idea]
Approximate each $\Omega_n$ by even step functions and use the lattice-time transfer-matrix construction: evenness yields a positive transfer matrix across $t=0$. Positivity is preserved under strong limits.
\end{proof}

\paragraph{Failure mode.} If $\Omega_n$ has a nontrivial odd part in $t$, the reflection symmetry is lost and OS positivity can fail.

\paragraph{Minimal model-preserving fixes.}
\begin{itemize}[leftmargin=*]
  \item \textbf{Symmetrization:} replace $T_{nn}(t)$ by $T^{(+)}_{nn}(t):=\tfrac12\big(T_{nn}(t)+T_{nn}(-t)\big)$. Then $\Omega_n(-t)=\Omega_n(t)$ and \eqref{eq:uniformLower} is preserved.
  \item \textbf{Stationarization:} replace $T_{nn}(t)$ by $T^{\mathrm{const}}_{nn}:=\inf_t T_{nn}(t)$, yielding a constant mass profile and automatic OS positivity.
\end{itemize}

\section{Safe generalization: off-diagonal mode couplings}\label{sec:offdiag}
Allow couplings via a kernel $K_{mn}(t-s)$ in the quadratic action
\[
S_E=\tfrac12\sum_{m,n}\iint \Phi_m(t)\,K_{mn}(t-s)\,\Phi_n(s)\,dt\,ds.
\]
A sufficient condition for OS positivity is: (i) \emph{evenness} $K_{mn}(t)=K_{mn}(-t)$, and (ii) \emph{positive type} (Bochner): the temporal Fourier transform $\widehat K(\xi)=[\widehat K_{mn}(\xi)]_{m,n}$ is positive semidefinite for almost every $\xi\in\R$. On the lattice, this corresponds to a block-Toeplitz, reflection-symmetric, positive semidefinite matrix.

\section{All-at-once theorem}\label{sec:allatonce}
\begin{theorem}[Uniform mass gap for the prime-indexed quadratic model]\label{thm:main}
Let $\Hcal=\bigoplus_{n}L^2(\R)$ and, for each $n$, define $H_n(t)$ by \eqref{eq:Hn} with $\Omega_n(t)^2$ given by \eqref{eq:Omega}. Assume \eqref{eq:uniformLower}. Then, for all $n,t$:
\[
\text{(i) } H_n(t)\text{ is essentially self-adjoint on }C_c^\infty,\quad
\text{(ii) } \inf\spec H_n(t)\ge \tfrac12\sqrt{\delta_n},\quad
\text{(iii) } \Delta_n(t)=\Omega_n(t)\ge \sqrt{\delta_n}.
\]
If $\delta_*:=\inf_n\delta_n>0$, then for the normal-ordered direct-sum dynamics on the symmetric Fock space, the many-body excitation gap satisfies $\Delta(t)\ge\sqrt{\delta_*}>0$ for all $t$. If, in addition, each $\Omega_n(t)$ is even in $t$, the associated Euclidean Gaussian measure is Osterwalder--Schrader reflection positive.
\end{theorem}

\section*{Acknowledgments}
We thank colleagues for discussions on operator domains, non-autonomous form methods, and reflection positivity. Any remaining errors are ours.

\appendix
\section{Reference facts (for quick citation)}\label{app:ref}
\begin{itemize}[leftmargin=*]
  \item \textbf{KLMN Theorem:} Closed form bounds imply self-adjoint operator bounds and spectral comparisons.
  \item \textbf{Faris--Lavine / Kato--Rellich:} Essential self-adjointness for $-\tfrac12\partial_x^2 + a x^2$ with $a>0$; compact resolvent for confining quadratics.
  \item \textbf{Kato (time-dependent forms):} Stability of time-dependent quadratic forms yields a unique propagator with time-independent form domain (see Kato; also Pazy, Chs.~4--5; Engel--Nagel).
  \item \textbf{OS positivity (Gaussian):} Even-in-time, positive-type kernels produce reflection positivity; lattice step-function approximations and transfer matrices provide a constructive proof.
\end{itemize}

% Bibliography placeholder (choose your system)
% \bibliographystyle{alpha}
% \bibliography{refs}

\end{document}
